\section{Conclusion}

Paired-Acquisition Neural Factorization is presented here as a complete computational pathology research pipeline extending from scanner-aware representations to whole-slide and multi-institutional learning. The representation stage uses same-tissue multi-scanner observations to route acquisition information between explicit tissue-oriented and acquisition-oriented branches. The whole-slide stage uses TransnnMIL and matched MIL baselines over a verified 10,611-slide PANDA feature pipeline. The institutional stage uses PathologyFL, dominance-aware aggregation, and center-weighting experiments to study how site influence changes when sample volume and task alignment diverge.

The strongest results are specific and complementary. On SCORPION, PA-NF has a registered controlled advantage over an equal-capacity two-branch neural control, reducing tissue-branch scanner balanced accuracy by 0.3108 while satisfying the predefined same-region retrieval noninferiority criterion and retaining scanner information in the acquisition branch. On PANDA, a fixed dominance-aware detector improves over FedAvg in specified dominant-site stress regimes and transfers without retuning from random label noise to systematic ordinal threshold shift. On CAMELYON17/WILDS, equal-client or dominant-source-adjusted weighting improves held-out-center accuracy over sample-proportional weighting in both tested feature families. TransnnMIL adds a stable custom whole-slide architecture, while the external canine and WSI-NCA experiments define important boundaries around what has and has not yet been established.

The resulting contribution is not a universal leaderboard or a clinical system. It is an end-to-end experimental framework for asking how information is represented and aggregated across scanners, tissue, slides, and institutions---and for replacing implicit assumptions about scanner invariance or sample-volume authority with explicit, testable neural and statistical mechanisms.
